\chapter{Stabilising the Finite-Strain Damage Route}
\label{ch:finite-strain-damage-stabilization}

\section{Motivation}

Chapter~\ref{ch:finite-strain-damage-relation} introduced the first real
finite-strain inelastic continuum law in the library, but intentionally in its
most conservative form: a scalar local problem and a degraded secant tangent.
That was the right first cut to validate the architecture, but it still left
two weaknesses:
\begin{itemize}
  \item the local problem was too special, and therefore did not exercise the
        vector-valued local-problem framework that future continuum
        damage-plasticity and fracture-oriented laws will need;
  \item the tangent was only secant-like, so active loading did not expose a
        consistent linearisation of the damage update.
\end{itemize}

This chapter documents the stabilised version of that route.  The objective is
not to over-specialise the model, but to keep the implementation open to
compile-time extension while removing the most obvious provisional pieces.

\section{State and Local Unknown}

The algorithmic state remains
\[
  \alpha = \{\kappa, d\},
\]
where $\kappa$ is the history variable and $d$ is the scalar damage variable.

The stabilised local problem no longer solves only for $\kappa$.  Instead, it
solves for the vector
\[
  \mathbf{x} =
  \begin{bmatrix}
    \kappa \\ d
  \end{bmatrix}.
\]

This is a small change mathematically, but a large one architecturally:
the finite-strain damage law now follows the same fixed-size vector local
problem interface that future coupled constitutive models will need.

\section{Vector Local Problem}

For the present rate-independent isotropic law,
\[
  \kappa_{n+1} = \max(\kappa_n, Y_{n+1}),
  \qquad
  d_{n+1} = D(\kappa_{n+1}).
\]

The local problem is written as
\[
  \mathbf{R}(\mathbf{x}) =
  \begin{bmatrix}
    R_1(\kappa,d) \\
    R_2(\kappa,d)
  \end{bmatrix}
  =
  \begin{bmatrix}
    \kappa - \max(\kappa_n, Y_{n+1}) \\
    d - D(\kappa)
  \end{bmatrix}
  = \mathbf{0}.
\]

Its Jacobian is
\[
  \mathbf{K} =
  \frac{\partial \mathbf{R}}{\partial \mathbf{x}}
  =
  \begin{bmatrix}
    1 & 0 \\
    -D'(\kappa) & 1
  \end{bmatrix}.
\]

In the current model this system is triangular, so Newton still converges in a
trivial manner.  However, that is no longer the architectural point.  The
important point is that the constitutive infrastructure now sees a
\emph{vector-valued local problem} rather than a special scalar shortcut.

\section{Consistent Tangent During Active Damage Loading}

The constitutive law is
\[
  \mathbf{S}(\mathbf{E},d) = g(d)\,\mathbf{S}_0(\mathbf{E}),
\]
with $\mathbf{S}_0 = \partial \psi_0 / \partial \mathbf{E}$ and
$g(d)$ the scalar degradation function.

When the loading is active, i.e. when the updated state satisfies
\[
  \kappa_{n+1} = Y_{n+1},
\]
the algorithmic tangent is no longer purely secant.  Differentiating
\[
  \mathbf{S} = g(d(\kappa(\mathbf{E})))\,\mathbf{S}_0(\mathbf{E}),
\]
gives
\[
  \mathbb{C}_{\text{alg}}
  =
  g(d)\,\mathbb{C}_0
  +
  \mathbf{S}_0 \otimes
  \frac{\partial g}{\partial \mathbf{E}}.
\]

Using the chain rule,
\[
  \frac{\partial g}{\partial \mathbf{E}}
  =
  \frac{\partial g}{\partial d}
  \frac{\partial d}{\partial \kappa}
  \frac{\partial \kappa}{\partial Y}
  \frac{\partial Y}{\partial \mathbf{E}}.
\]

For the present rate-independent model, during active loading
$\partial \kappa / \partial Y = 1$, so
\[
  \frac{\partial g}{\partial \mathbf{E}}
  =
  g'(d)\,D'(\kappa)\,\frac{\partial Y}{\partial \mathbf{E}}.
\]

During unloading or below threshold, the history variable is frozen and the
additional term vanishes.  The tangent then reduces again to the degraded
secant form
\[
  \mathbb{C}_{\text{alg}} = g(d)\,\mathbb{C}_0.
\]

\section{Driving-Force Gradient}

The current driving variable is the norm of the positive spectral part of the
Green--Lagrange strain,
\[
  Y(\mathbf{E}) = \left\lVert \langle \mathbf{E} \rangle_+ \right\rVert.
\]

Let
\[
  \mathbf{E}_+ = \sum_a \langle \lambda_a \rangle_+ \,\mathbf{n}_a \otimes
  \mathbf{n}_a,
\]
where $\lambda_a$ and $\mathbf{n}_a$ are the principal strains and principal
directions of $\mathbf{E}$.  For $Y > 0$,
\[
  \frac{\partial Y}{\partial \mathbf{E}}
  =
  \frac{\mathbf{E}_+}{\lVert \mathbf{E}_+ \rVert}.
\]

That is the derivative now implemented in
\texttt{PositiveGreenLagrangeEquivalentStrain}.  In code, the derivative is
exposed with respect to the engineering-Voigt strain vector used by the
element formulation, which keeps the constitutive tangent directly compatible
with the existing assembly path.

\section{Compile-Time Extension Points}

The stabilised implementation preserves the library's compile-time openness:
\begin{itemize}
  \item \textbf{Driving force policy}: any policy satisfying
        \texttt{DamageDrivingForcePolicy} may still drive the law.
  \item \textbf{Differentiable driving force}: if a policy additionally
        exposes \texttt{gradient(...)} or
        \texttt{gradient\_wrt\_engineering(...)}, the law can build a more
        consistent tangent.  If not, it falls back cleanly.
  \item \textbf{Damage evolution policy}: any
        \texttt{DamageEvolutionPolicy} remains admissible.
  \item \textbf{Differentiable evolution}: if the policy exposes
        \texttt{damage\_derivative(...)} and
        \texttt{degradation\_derivative(...)}, the rank-one correction is
        activated during loading.  Otherwise the law degrades to the secant
        tangent without breaking the interface.
  \item \textbf{Local solver}: Newton is still only one admissible local
        algorithm.  The tests continue to inject an exact direct solver to
        verify that the local-problem interface remains solver-agnostic.
\end{itemize}

This is the key architectural point: the law does not become less general by
becoming more consistent.  It becomes more faithful \emph{when the injected
policies provide the required derivatives}, and remains conservative otherwise.

\section{Verification Performed}

The regression suite now verifies:
\begin{itemize}
  \item the vector local problem produces the same state, response, and
        tangent under both Newton and an injected exact solver;
  \item active loading produces an algorithmic tangent that matches finite
        differences of the fully integrated constitutive response;
  \item unloading reduces to the degraded secant tangent, as expected from a
        frozen history variable;
  \item the spatial updated-Lagrangian path remains correct for stress and
        tangent under unloading, where the secant form is exact for the
        present model.
\end{itemize}

\section{What Remains Open}

Even after this stabilisation, the route is intentionally not closed:
\begin{itemize}
  \item the model remains local and scalar;
  \item no crack-band or nonlocal regularisation is introduced yet;
  \item the spatial tangent is still obtained through push-forward of the
        reference tangent, which is acceptable for the present path but not yet
        the final word for general finite-strain inelasticity;
  \item coupled finite-strain plasticity-damage still requires a higher-rank
        local state and a genuinely vector-valued constitutive residual.
\end{itemize}

That said, the current route is now stable enough to be a credible base layer
for those next developments.

\section{References}

\begin{thebibliography}{9}

\bibitem{simo1987damage2}
J.~C. Simo and J.~W. Ju,
\emph{Strain- and stress-based continuum damage models. I. Formulation},
International Journal of Solids and Structures, 23(7), 821--840, 1987.

\bibitem{miehe1995damage2}
C.~Miehe,
\emph{Discontinuous and continuous damage evolution in Ogden-type large-strain
elastic materials},
European Journal of Mechanics A/Solids, 14(4), 697--720, 1995.

\bibitem{simohughes1998damage2}
J.~C. Simo and T.~J.~R. Hughes,
\emph{Computational Inelasticity},
Springer, 1998.

\bibitem{bonetwood2008damage2}
J.~Bonet, R.~D. Wood, and A.~J. Gil,
\emph{Nonlinear Solid Mechanics for Finite Element Analysis: Statics},
2nd ed., Cambridge University Press, 2008.

\bibitem{desouzaneto2008damage2}
E.~A. de Souza Neto, D.~Peri\'{c}, and D.~R.~J. Owen,
\emph{Computational Methods for Plasticity: Theory and Applications},
Wiley, 2008.

\end{thebibliography}
